\documentclass[12pt]{article}

\usepackage{fancyhdr}

\pagestyle{fancy}
\fancyhead[l]{DispatchNet}
\fancyhead[c]{Project TRE}
\fancyhead[r]{\today}
\fancyfoot[c]{\thepage}

\renewcommand\bf[1]{\textbf{#1}}
\renewcommand\it[1]{\textit{#1}}
\newcommand\un[1]{\underline{#1}}

\begin {document}
\tableofcontents

\newpage
\section{Objectives}
	\subsection{Rail System Modeling}
		The system aims to model a real railway layout in terms of:
		\begin{itemize}
			\item \bf{Lines.}
				Collection of tracks that run without interference between two \bf{Junctions} or \bf{Stations}
			\item \bf{Junctions.} Anonymous \bf{Stations} where lines connect
			\item \bf{Stations.} Places for trains to stop and load passengers or cargo	
		\end{itemize}
		
		\it{Boundary:} 
		Initially, we are not aiming to model each individual section of track, instead we abstract them to \bf{Lines} and \bf{Stations} to group them, reducing complexity.
		In the future, the system may be expanded to allow for finer granularity
		
	\subsection{Travel Information}
		The system must allow anyone to view relevant travel information on any relevant \un{Passenger} journey on the network

	\subsection{Service Planning}
		The system must allow \un{Companies} to register services to be ran along the network

	\subsection{Train Definitions}
		For ease of use, we must keep a record of existing train types for use in services in order to calculate capacity and enforce other physical constraints

\newpage
\section{Non-Functional Requirements}
	\subsection{Topology}
		Following \bf{1.1}, the system must be able to store the layout of the railway network as a directed mesh.
		Each node in the mesh should not have restrictions over the number of incoming and outgoing connections, beyond having at least one of each (not being orphaned or a dead-end)
	\subsection{Unique IDs}
		Following \bf{1.1} and in reference to \bf{2.1}, the system should store each element in the network with unique IDs for each type of element.
	
\section{Functional Requirements}
	\subsection{Creating a Service}
		Following \bf{1.3} and in reference to \bf{2.1}, the system must enavle the easy creation of a route between different stations
	\subsection{Creating a Train}
		The system must allow \un{Companies} to register a train type.
		Upon definition, a train type should include the following fields:
		\begin{itemize}
			\item The train series identificator
			\item (Optionally) The commercial name of the train
			\item The set of service-types that the train can be used for
			\item The lenght
			\item The maximum operating speed of a train
			\item Whether it is a \un{Passenger} or \un{Freight} train
			\item The capacities, excluding staff, of people or cargo respective to the previous train type
		\end{itemize}
\end {document}
